\documentclass[12pt,a4paper]{report}

% ---------------------------------------------------------------------
% Packages
% ---------------------------------------------------------------------
\usepackage[utf8]{inputenc}
\usepackage[T1]{fontenc}
\usepackage[french]{babel}
\usepackage[margin=2.5cm,headheight=14pt]{geometry}
\usepackage{graphicx}
\usepackage{tikz}
\usetikzlibrary{positioning, arrows.meta, shapes.geometric}
\usepackage{xcolor}
\usepackage{booktabs}
\usepackage{enumitem}
\usepackage{titlesec}
\usepackage{fancyhdr}
\usepackage{hyperref}
\usepackage{caption}

% ---------------------------------------------------------------------
% Couleurs
% ---------------------------------------------------------------------
\definecolor{accent}{HTML}{1F4E79}
\definecolor{accentlight}{HTML}{5B9BD5}
\definecolor{textgray}{HTML}{666666}

% ---------------------------------------------------------------------
% Mise en forme des titres (sobre, sans cadres)
% ---------------------------------------------------------------------
\titleformat{\chapter}[hang]{\normalfont\Huge\bfseries\color{accent}}{\thechapter.}{20pt}{\Huge}
\titleformat{\section}{\normalfont\Large\bfseries\color{accent}}{\thesection}{1em}{}
\titleformat{\subsection}{\normalfont\large\bfseries\color{accentlight}}{\thesubsection}{1em}{}

% ---------------------------------------------------------------------
% En-têtes / pieds de page
% ---------------------------------------------------------------------
\pagestyle{fancy}
\fancyhf{}
\renewcommand{\headrulewidth}{0.4pt}
\renewcommand{\footrulewidth}{0pt}
\fancyhead[L]{\small NeurAR -- Rapport de projet}
\fancyhead[R]{\small \leftmark}
\fancyfoot[C]{\thepage}

% ---------------------------------------------------------------------
% Hyperliens
% ---------------------------------------------------------------------
\hypersetup{
    colorlinks=true,
    linkcolor=accent,
    citecolor=accent,
    urlcolor=accent,
    pdftitle={NeurAR - Rapport de projet},
    pdfauthor={CHAIBOU SAIDOU ABDOULAYE}
}

\begin{document}

% =======================================================================
% Page de titre
% =======================================================================
\begin{titlepage}
    \centering
    \vspace*{1.5cm}

    {\Large \textsc{Projet de fin d'\'etudes}}\\[0.3cm]
    {\small R\'ealit\'e augment\'ee \& Intelligence artificielle}\\[2.5cm]

    {\Huge \bfseries \color{accent} NeurAR}\\[0.6cm]
    {\Large Visualisation en r\'ealit\'e augment\'ee\\ d'un r\'eseau de neurones en fonctionnement}\\[2.5cm]

    \begin{center}
    \begin{tabular}{ll}
        \textbf{R\'ealis\'e par :} & CHAIBOU SAIDOU ABDOULAYE \\[0.4cm]
        \textbf{Encadr\'e par :} & Pr.~Tawfik Masrour \\
                                  & Pr.~Ajebli Aymane \\
                                  & Pr.~Maalouf Imad \\
    \end{tabular}
    \end{center}

    \vfill

    {\large Ann\'ee universitaire 2025 -- 2026}

\end{titlepage}

% =======================================================================
% Résumé
% =======================================================================
\chapter*{R\'esum\'e}
\addcontentsline{toc}{chapter}{R\'esum\'e}

Les r\'eseaux de neurones sont aujourd'hui au c\oe ur de la plupart des
applications d'intelligence artificielle, mais leur fonctionnement interne
reste difficile \`a appr\'ehender pour les \'etudiants et les personnes
d\'ecouvrant le domaine. Le projet \textbf{NeurAR} propose une solution
originale \`a ce probl\`eme : visualiser, en \textbf{r\'ealit\'e augment\'ee},
un r\'eseau de neurones de type \textit{Multi-Layer Perceptron} (MLP) en train
de fonctionner, couche par couche, en temps r\'eel.

Le syst\`eme repose sur trois composants principaux : un \textbf{mod\`ele
PyTorch} entra\^in\'e sur le jeu de donn\'ees MNIST, un \textbf{backend
FastAPI} exposant une API REST et un canal WebSocket temps r\'eel, et une
\textbf{application Unity} en r\'ealit\'e augment\'ee qui re\c{c}oit ces
donn\'ees et les transforme en une repr\'esentation 3D anim\'ee et
interactive. L'utilisateur peut ainsi observer en direct la propagation
avant (\textit{forward pass}) et la r\'etropropagation (\textit{backward
pass}) du r\'eseau, manipuler la sc\`ene en AR (rotation, zoom, s\'election
de neurone) et suivre l'\'evolution de la perte et de la pr\'ecision du
mod\`ele.

Ce rapport pr\'esente le contexte et les objectifs du projet, l'architecture
technique mise en place, les choix d'impl\'ementation pour le mod\`ele de
r\'eseau de neurones, le protocole de communication temps r\'eel, ainsi que
l'application Unity AR. Il revient \'egalement sur les principales
difficult\'es rencontr\'ees, les solutions apport\'ees, et propose des pistes
d'am\'elioration pour la suite du projet.

% =======================================================================
% Table des matières
% =======================================================================
\tableofcontents

% =======================================================================
\chapter{Introduction}

\section{Contexte}

L'intelligence artificielle, et en particulier l'apprentissage profond
(\textit{deep learning}), occupe une place de plus en plus importante dans
de nombreux domaines : reconnaissance d'images, traitement du langage,
syst\`emes de recommandation, etc. Cependant, les r\'eseaux de neurones sont
souvent per\c{c}us comme des \textbf{bo\^ites noires} : on conna\^it leurs
entr\'ees et leurs sorties, mais leur fonctionnement interne (propagation des
activations, mise \`a jour des poids) reste abstrait, surtout pour les
personnes qui d\'ecouvrent le domaine.

Par ailleurs, la \textbf{r\'ealit\'e augment\'ee (RA)} permet de superposer
des objets virtuels au monde r\'eel et offre de nouvelles possibilit\'es
p\'edagogiques : manipulation d'objets 3D, immersion, interaction naturelle.
Combiner ces deux univers -- intelligence artificielle et r\'ealit\'e
augment\'ee -- ouvre la voie \`a des outils p\'edagogiques innovants pour
mieux comprendre le fonctionnement des mod\`eles d'apprentissage automatique.

\section{Probl\'ematique}

Comment permettre \`a un utilisateur de \textbf{visualiser et de
manipuler}, en temps r\'eel et en trois dimensions, le fonctionnement
interne d'un r\'eseau de neurones (propagation des activations, poids,
biais, m\'etriques d'apprentissage), tout en b\'en\'eficiant de
l'immersion offerte par la r\'ealit\'e augment\'ee ?

\section{Objectifs du projet}

Le projet NeurAR poursuit les objectifs suivants :

\begin{itemize}
    \item Concevoir et entra\^iner un r\'eseau de neurones de type MLP sur le
    jeu de donn\'ees MNIST \`a l'aide de PyTorch ;
    \item D\'evelopper un backend exposant l'\'etat complet du r\'eseau
    (architecture, activations, poids, biais, m\'etriques) via une API REST
    et un canal WebSocket temps r\'eel ;
    \item Concevoir une application Unity en r\'ealit\'e augment\'ee
    capable de placer le r\'eseau de neurones dans l'environnement r\'eel et
    de le repr\'esenter en 3D ;
    \item Animer visuellement la propagation avant (\textit{forward pass})
    et la r\'etropropagation (\textit{backward pass}), couche par couche ;
    \item Permettre \`a l'utilisateur d'interagir avec la sc\`ene
    (rotation, zoom, s\'election d'un neurone) ;
    \item Garantir une communication fluide, synchronis\'ee et fiable entre
    le backend Python et le client Unity.
\end{itemize}

\section{Plan du rapport}

Ce rapport est organis\'e comme suit. Le chapitre~\ref{chap:archi} pr\'esente
l'architecture g\'en\'erale de la solution. Le chapitre~\ref{chap:backend}
d\'ecrit le mod\`ele de r\'eseau de neurones et le backend FastAPI. Le
chapitre~\ref{chap:comm} d\'etaille le protocole de communication temps
r\'eel. Le chapitre~\ref{chap:unity} pr\'esente l'application Unity en
r\'ealit\'e augment\'ee. Le chapitre~\ref{chap:resultats} pr\'esente les
fonctionnalit\'es r\'ealis\'ees et les apports du projet. Le
chapitre~\ref{chap:difficultes} revient sur les difficult\'es rencontr\'ees
et les solutions apport\'ees. Enfin, le chapitre~\ref{chap:conclusion}
conclut ce rapport et propose des perspectives d'am\'elioration.


% =======================================================================
\chapter{Architecture g\'en\'erale de la solution}
\label{chap:archi}

\section{Vue d'ensemble}

NeurAR repose sur une architecture \textbf{client-serveur} compos\'ee de
trois grands blocs : le \textbf{mod\`ele PyTorch}, le \textbf{backend
FastAPI} et le \textbf{client Unity} en r\'ealit\'e augment\'ee. La
figure~\ref{fig:archi-globale} sch\'ematise les \'echanges entre ces
composants.

\begin{figure}[h!]
\centering
\begin{tikzpicture}[
    node distance=1.6cm and 2.2cm,
    box/.style={rectangle, rounded corners, draw=accent, fill=accent!10,
                 minimum width=3.6cm, minimum height=1.6cm, align=center,
                 font=\small},
    arrow/.style={-Stealth, thick, accent}
]
    \node[box] (unity) {Application Unity\\(client AR)\\rendu 3D + interactions};
    \node[box, right=of unity] (backend) {Backend FastAPI\\API REST + WebSocket\\(\texttt{/ws})};
    \node[box, right=of backend] (model) {Mod\`ele PyTorch\\MLP (NeuroARNet)};

    \draw[arrow] ([yshift=2mm]unity.east) -- node[above, font=\scriptsize]{commandes (forward, backward, reset)} ([yshift=2mm]backend.west);
    \draw[arrow] ([yshift=-2mm]backend.west) -- node[below, font=\scriptsize]{\'etat du r\'eseau (JSON)} ([yshift=-2mm]unity.east);

    \draw[arrow] ([yshift=2mm]backend.east) -- node[above, font=\scriptsize]{appels internes} ([yshift=2mm]model.west);
    \draw[arrow] ([yshift=-2mm]model.west) -- node[below, font=\scriptsize]{activations, poids, biais} ([yshift=-2mm]backend.east);
\end{tikzpicture}
\caption{Architecture g\'en\'erale de NeurAR}
\label{fig:archi-globale}
\end{figure}

\section{R\^ole de chaque composant}

\subsection{Le mod\`ele PyTorch}
Le mod\`ele est un \textit{Multi-Layer Perceptron} (MLP) entra\^in\'e sur le
jeu de donn\'ees MNIST. Il est responsable du calcul des activations, des
poids et des biais de chaque couche, ainsi que des m\'etriques
d'apprentissage (perte, pr\'ecision).

\subsection{Le backend FastAPI}
Le backend joue le r\^ole d'interm\'ediaire entre le mod\`ele et le client.
Il expose :
\begin{itemize}
    \item une \textbf{API REST} pour les requ\^etes ponctuelles
    (sant\'e du serveur, architecture du mod\`ele, instantan\'e du r\'eseau,
    forward pass) ;
    \item un \textbf{canal WebSocket} (\texttt{/ws}) pour les interactions
    temps r\'eel (lancement d'un forward pass ou backward pass, r\'einitialisation
    du mod\`ele, diffusion de l'\'etat du r\'eseau couche par couche).
\end{itemize}

\subsection{L'application Unity (client AR)}
L'application Unity re\c{c}oit les donn\'ees JSON envoy\'ees par le backend,
les interpr\`ete et construit une repr\'esentation 3D du r\'eseau de
neurones plac\'ee dans l'environnement r\'eel via la r\'ealit\'e augment\'ee.
Elle g\`ere \'egalement les interactions de l'utilisateur (rotation, zoom,
s\'election de neurone).

\section{Technologies utilis\'ees}

\begin{table}[h!]
\centering
\begin{tabular}{ll}
\toprule
\textbf{Composant} & \textbf{Technologies} \\
\midrule
Mod\`ele de r\'eseau de neurones & Python, PyTorch, torchvision, MNIST \\
Backend & Python, FastAPI, WebSocket, Pydantic \\
Client AR & Unity 6000.4.3f1, AR Foundation, C\# \\
Format d'\'echange & JSON (camelCase, compatible C\#) \\
\bottomrule
\end{tabular}
\caption{Technologies utilis\'ees dans le projet NeurAR}
\end{table}


% =======================================================================
\chapter{Le mod\`ele de r\'eseau de neurones et le backend}
\label{chap:backend}

\section{Architecture du r\'eseau (NeuroARNet)}

Le r\'eseau de neurones impl\'ement\'e, nomm\'e \texttt{NeuroARNet}, est un
\textbf{MLP enti\`erement connect\'e et configurable}. Son architecture est
d\'efinie par une simple liste d'entiers (par exemple
\texttt{[784, 128, 64, 10]}), modifiable via le fichier de configuration
\texttt{.env}. Pour cette architecture par d\'efaut :

\begin{itemize}
    \item \textbf{Couche d'entr\'ee} : 784 neurones (image MNIST $28\times28$
    aplatie) ;
    \item \textbf{Couche cach\'ee 1} : \texttt{Linear(784 $\rightarrow$ 128)}
    + \texttt{BatchNorm} + \texttt{ReLU} ;
    \item \textbf{Couche cach\'ee 2} : \texttt{Linear(128 $\rightarrow$ 64)}
    + \texttt{BatchNorm} + \texttt{ReLU} ;
    \item \textbf{Couche de sortie} : \texttt{Linear(64 $\rightarrow$ 10)} +
    \texttt{Softmax} (probabilit\'es des 10 chiffres).
\end{itemize}

La figure~\ref{fig:mlp} illustre cette architecture sous forme simplifi\'ee.

\begin{figure}[h!]
\centering
\begin{tikzpicture}[
    neuron/.style={circle, draw=accent, fill=accent!20, minimum size=0.5cm},
    neuronout/.style={circle, draw=accentlight, fill=accentlight!30, minimum size=0.5cm},
    layerlabel/.style={font=\small, align=center}
]
    % Couche d'entrée (4 neurones représentés)
    \foreach \i in {1,...,4}
        \node[neuron] (I-\i) at (0, -\i*0.8) {};
    \node[layerlabel] at (0, -4.6) {Entr\'ee\\(784)};

    % Couche cachée 1 (5 neurones représentés)
    \foreach \i in {1,...,5}
        \node[neuron] (H1-\i) at (3, -\i*0.8+0.4) {};
    \node[layerlabel] at (3, -4.6) {Cach\'ee 1\\(128)\\Linear+BN+ReLU};

    % Couche cachée 2 (4 neurones représentés)
    \foreach \i in {1,...,4}
        \node[neuron] (H2-\i) at (6, -\i*0.8) {};
    \node[layerlabel] at (6, -4.6) {Cach\'ee 2\\(64)\\Linear+BN+ReLU};

    % Couche de sortie (3 neurones représentés)
    \foreach \i in {1,...,3}
        \node[neuronout] (O-\i) at (9, -\i*0.8-0.4) {};
    \node[layerlabel] at (9, -4.6) {Sortie\\(10)\\Linear+Softmax};

    % Connexions entrée -> cachée 1
    \foreach \i in {1,...,4}
        \foreach \j in {1,...,5}
            \draw[gray!40, thin] (I-\i) -- (H1-\j);

    % Connexions cachée 1 -> cachée 2
    \foreach \i in {1,...,5}
        \foreach \j in {1,...,4}
            \draw[gray!40, thin] (H1-\i) -- (H2-\j);

    % Connexions cachée 2 -> sortie
    \foreach \i in {1,...,4}
        \foreach \j in {1,...,3}
            \draw[gray!40, thin] (H2-\i) -- (O-\j);
\end{tikzpicture}
\caption{Repr\'esentation simplifi\'ee de l'architecture du MLP NeuroARNet}
\label{fig:mlp}
\end{figure}

\section{Suivi des activations, poids et biais}

\`A chaque \textit{forward pass}, le mod\`ele conserve les activations de
toutes les couches (y compris la couche d'entr\'ee). Le
\texttt{ModelManager} expose ensuite ces informations sous une forme
exploitable par Unity :

\begin{itemize}
    \item les \textbf{activations} de chaque couche sont normalis\'ees entre
    0 et 1 (min-max scaling) afin d'\^etre directement utilis\'ees pour
    l'intensit\'e visuelle des neurones ;
    \item les \textbf{poids} de chaque couche sont normalis\'es entre $-1$ et
    $1$ via la fonction $\tanh$, ce qui conserve leur signe tout en bornant
    leur amplitude pour l'affichage des connexions ;
    \item les \textbf{biais} de chaque neurone sont transmis tels quels.
\end{itemize}

\section{Entra\^inement et m\'etriques}

Le mod\`ele est entra\^in\'e/optimis\'e sur des \'echantillons du jeu de
donn\'ees \textbf{MNIST}, t\'el\'echarg\'e automatiquement au premier
lancement. \`A chaque \textit{backward pass}, un optimiseur SGD met \`a jour
les poids du r\'eseau \`a partir d'un \'echantillon, et les m\'etriques
suivantes sont recalcul\'ees :

\begin{itemize}
    \item la \textbf{perte} (\textit{loss}), calcul\'ee via la
    \textit{negative log-likelihood} ;
    \item la \textbf{pr\'ecision} (\textit{accuracy}), \'egale \`a 1 si la
    classe pr\'edite correspond \`a l'\'etiquette attendue, 0 sinon.
\end{itemize}

\section{Mod\'elisation des donn\'ees \'echang\'ees}

Pour garantir la compatibilit\'e directe avec les classes C\# du c\^ot\'e
Unity, les mod\`eles de donn\'ees Python (Pydantic) g\'en\`erent
automatiquement des alias en \texttt{camelCase}. Trois structures
principales sont d\'efinies :

\begin{itemize}
    \item \texttt{NeuronData} : un neurone (couche, index, activation, biais,
    fonction d'activation) ;
    \item \texttt{LayerData} : une couche compl\`ete (type, liste de
    neurones, poids) ;
    \item \texttt{NetworkData} : l'\'etat complet du r\'eseau (nom du
    mod\`ele, nombre de couches, liste des couches, phase, perte,
    pr\'ecision).
\end{itemize}

\section{API REST}

Le tableau~\ref{tab:rest} r\'esume les points d'entr\'ee REST expos\'es par
le backend.

\begin{table}[h!]
\centering
\begin{tabular}{lll}
\toprule
\textbf{M\'ethode} & \textbf{Route} & \textbf{Description} \\
\midrule
GET & \texttt{/health} & V\'erifie que le serveur et le mod\`ele sont op\'erationnels \\
GET & \texttt{/network} & Retourne un instantan\'e complet du r\'eseau \\
POST & \texttt{/forward} & Ex\'ecute un forward pass sur une entr\'ee donn\'ee \\
GET & \texttt{/architecture} & Retourne la configuration et le nombre de param\`etres \\
\bottomrule
\end{tabular}
\caption{Points d'entr\'ee de l'API REST}
\label{tab:rest}
\end{table}


% =======================================================================
\chapter{Communication temps r\'eel}
\label{chap:comm}

\section{Le canal WebSocket}

La communication temps r\'eel entre le backend et le client Unity repose sur
un canal \textbf{WebSocket} (\texttt{/ws}). Le client envoie un message JSON
contenant un champ \texttt{action}, et le serveur r\'epond par un ou
plusieurs messages \texttt{NetworkData}.

\begin{table}[h!]
\centering
\begin{tabular}{lp{9cm}}
\toprule
\textbf{Action} & \textbf{Comportement} \\
\midrule
\texttt{get\_network} & Envoie l'\'etat complet du r\'eseau (phase \texttt{init}) \\
\texttt{forward\_pass} & Ex\'ecute un forward pass sur un \'echantillon MNIST
al\'eatoire et diffuse le r\'esultat \textbf{couche par couche} (d\'elai de
100~ms entre chaque couche), puis l'\'etat final complet \\
\texttt{backward\_pass} & Ex\'ecute un backward pass (mise \`a jour des
poids) et diffuse le nouvel \'etat \\
\texttt{reset} & R\'einitialise le mod\`ele (nouveaux poids al\'eatoires) et
diffuse le nouvel \'etat \\
\bottomrule
\end{tabular}
\caption{Actions support\'ees par le canal WebSocket}
\end{table}

\section{Sc\'enario d'animation du forward pass}

La figure~\ref{fig:sequence} pr\'esente le sc\'enario d'\'echange lors d'un
\textit{forward pass}. L'envoi couche par couche, avec un l\'eger d\'elai,
permet \`a Unity d'animer progressivement la propagation de l'activation \`a
travers le r\'eseau.

\begin{figure}[h!]
\centering
\begin{tikzpicture}[
    node distance=1.4cm,
    every node/.style={font=\small},
    actor/.style={rectangle, rounded corners, draw=accent, fill=accent!10,
                   minimum width=3cm, minimum height=0.9cm, align=center},
    arrow/.style={-Stealth, thick, accent}
]
    \node[actor] (u) at (0,0) {Unity (client)};
    \node[actor] (b) at (8,0) {Backend FastAPI};

    \draw[arrow] (u.east) -- node[above]{\{action: forward\_pass\}} (b.west);

    \draw[dashed, gray] (0,-1) -- (0,-6.5);
    \draw[dashed, gray] (8,-1) -- (8,-6.5);

    \draw[arrow] (8,-2) -- node[above]{couche 0 (entr\'ee)} (0,-2);
    \draw[arrow] (8,-3) -- node[above]{couche 1 (apr\`es 100~ms)} (0,-3);
    \draw[arrow] (8,-4) -- node[above]{couche 2 (apr\`es 100~ms)} (0,-4);
    \draw[arrow] (8,-5) -- node[above]{couche 3 / sortie (apr\`es 100~ms)} (0,-5);
    \draw[arrow] (8,-6) -- node[above]{\'etat final complet (loss, accuracy)} (0,-6);
\end{tikzpicture}
\caption{S\'equence d'\'echange lors d'un forward pass}
\label{fig:sequence}
\end{figure}

\section{Gestion de la taille des messages}

Le snapshot complet du r\'eseau (poids inclus) repr\'esente un volume
important : pour une couche $784 \rightarrow 128$, cela correspond \`a
environ 100~000 valeurs, soit environ \textbf{2,4~Mo} en JSON. Cette
contrainte a un impact direct sur la configuration du client WebSocket
Unity (cf.~chapitre~\ref{chap:difficultes}).


% =======================================================================
\chapter{Application Unity en r\'ealit\'e augment\'ee}
\label{chap:unity}

\section{Vue d'ensemble}

L'application client est d\'evelopp\'ee avec \textbf{Unity 6000.4.3f1} et le
framework \textbf{AR Foundation}. Elle est organis\'ee en plusieurs modules
sp\'ecialis\'es, list\'es dans le tableau~\ref{tab:unity-scripts}.

\begin{table}[h!]
\centering
\begin{tabular}{lp{9.5cm}}
\toprule
\textbf{Module} & \textbf{R\^ole} \\
\midrule
\texttt{ARPlacementManager} & Placement du r\'eseau de neurones virtuel dans
l'environnement r\'eel \\
\texttt{ARInteractionManager} & Gestion des interactions utilisateur :
rotation, zoom, s\'election d'un neurone \\
\texttt{WebSocketClient} & Connexion et communication avec le backend via
WebSocket \\
\texttt{NetworkDataParser} & D\'es\'erialisation des messages JSON re\c{c}us
en objets \texttt{NetworkData} \\
\texttt{UnityMainThreadDispatcher} & Ex\'ecution des mises \`a jour re\c{c}ues
de mani\`ere asynchrone sur le thread principal Unity \\
\texttt{NetworkRenderer} & Construction et mise \`a jour de la sc\`ene 3D du
r\'eseau \\
\texttt{NetworkLayer} & Repr\'esentation d'une couche du r\'eseau \\
\texttt{NeuronVisual} & Repr\'esentation visuelle d'un neurone (couleur,
intensit\'e selon l'activation) \\
\texttt{ConnectionVisual} & Repr\'esentation visuelle d'une connexion
(\'epaisseur/couleur selon le poids) \\
\bottomrule
\end{tabular}
\caption{Principaux modules de l'application Unity}
\label{tab:unity-scripts}
\end{table}

\section{Cha\^ine de traitement c\^ot\'e Unity}

\begin{enumerate}
    \item Le \texttt{WebSocketClient} se connecte au backend et envoie les
    actions demand\'ees (\texttt{get\_network}, \texttt{forward\_pass}, etc.) ;
    \item Les messages JSON re\c{c}us sont transmis au
    \texttt{NetworkDataParser}, qui les convertit en objets
    \texttt{NetworkData} / \texttt{LayerData} / \texttt{NeuronData} ;
    \item L'\texttt{UnityMainThreadDispatcher} planifie la mise \`a jour de la
    sc\`ene sur le thread principal ;
    \item Le \texttt{NetworkRenderer} met \`a jour les \texttt{NeuronVisual}
    et \texttt{ConnectionVisual} en fonction des nouvelles activations,
    poids et biais ;
    \item L'\texttt{ARPlacementManager} et l'\texttt{ARInteractionManager}
    permettent \`a l'utilisateur de positionner et de manipuler la sc\`ene
    dans son environnement r\'eel.
\end{enumerate}

\section{Repr\'esentation visuelle}

Chaque \textbf{neurone} est repr\'esent\'e par un objet 3D dont la couleur ou
l'intensit\'e lumineuse refl\`ete sa valeur d'activation normalis\'ee.
Chaque \textbf{connexion} entre deux neurones est repr\'esent\'ee par un
segment dont l'apparence (couleur, \'epaisseur) refl\`ete le poids associ\'e
(normalis\'e entre $-1$ et $1$). Lors d'un forward pass, les couches
s'illuminent successivement, donnant une animation visuelle de la
propagation de l'information \`a travers le r\'eseau.


% =======================================================================
\chapter{Fonctionnalit\'es r\'ealis\'ees et apports du projet}
\label{chap:resultats}

\section{Fonctionnalit\'es r\'ealis\'ees}

\begin{itemize}
    \item Visualisation du r\'eseau de neurones, couche par couche, en 3D et
    en r\'ealit\'e augment\'ee ;
    \item Animation du \textbf{forward pass} : propagation visuelle de
    l'activation \`a travers les couches ;
    \item Animation du \textbf{backward pass} : mise \`a jour visuelle des
    poids apr\`es apprentissage ;
    \item Affichage en direct des valeurs d'activation, des poids et des
    biais de chaque neurone ;
    \item Suivi en temps r\'eel de la \textbf{perte (loss)} et de la
    \textbf{pr\'ecision (accuracy)} du mod\`ele ;
    \item R\'einitialisation du mod\`ele \`a la demande (nouveaux poids
    al\'eatoires) ;
    \item Interactions en r\'ealit\'e augment\'ee : placement de la sc\`ene,
    rotation, zoom et s\'election d'un neurone pour afficher ses d\'etails.
\end{itemize}

\section{Apports et points forts du projet}

\begin{itemize}
    \item \textbf{Approche p\'edagogique innovante} : le projet rend visible
    et manipulable le fonctionnement interne d'un r\'eseau de neurones,
    habituellement abstrait ;
    \item \textbf{Architecture d\'ecoupl\'ee et r\'eutilisable} : le backend
    FastAPI est ind\'ependant du client et pourrait servir d'autres
    applications (web, mobile) sans modification ;
    \item \textbf{Communication temps r\'eel performante} gr\^ace au
    WebSocket et \`a la diffusion couche par couche ;
    \item \textbf{Mod\`ele configurable} : l'architecture du MLP est
    modifiable simplement via le fichier \texttt{.env}, sans modification du
    code ;
    \item \textbf{Compatibilit\'e stricte avec Unity/C\#} gr\^ace \`a la
    g\'en\'eration automatique d'alias \texttt{camelCase} dans les mod\`eles
    Pydantic ;
    \item \textbf{Exp\'erience immersive} en r\'ealit\'e augment\'ee,
    utilisable sur smartphone ou casque AR compatible ;
    \item \textbf{Projet transversal} couvrant trois domaines compl\'ementaires :
    intelligence artificielle (PyTorch), d\'eveloppement backend (FastAPI) et
    applications immersives (Unity / AR Foundation).
\end{itemize}


% =======================================================================
\chapter{Difficult\'es rencontr\'ees et solutions apport\'ees}
\label{chap:difficultes}

\begin{itemize}
    \item \textbf{Volume de donn\'ees important} : un instantan\'e complet du
    r\'eseau (poids inclus) repr\'esente environ 2,4~Mo en JSON.\\
    \textit{Solution} : augmentation de la limite de taille des messages
    accept\'es par le client WebSocket Unity (\texttt{NativeWebSocket} /
    \texttt{websocket-sharp}) pour \'eviter toute d\'econnexion lors de la
    r\'eception du r\'eseau complet.

    \item \textbf{Synchronisation entre le calcul backend et l'animation
    Unity} : un envoi instantan\'e de toutes les couches rendrait
    l'animation illisible.\\
    \textit{Solution} : diffusion des couches une \`a une, avec un d\'elai
    contr\^ol\'e de 100~ms entre chaque couche, pour produire une animation
    fluide et compr\'ehensible.

    \item \textbf{Alignement strict entre les mod\`eles Python et les classes
    C\# Unity} : toute diff\'erence de nommage entre les champs JSON et les
    propri\'et\'es C\# emp\^eche la d\'es\'erialisation correcte.\\
    \textit{Solution} : mise en place d'une classe de base
    \texttt{CamelModel} g\'en\'erant automatiquement des alias
    \texttt{camelCase} pour chaque champ, garantissant une correspondance
    exacte avec \texttt{NetworkData}, \texttt{LayerData} et
    \texttt{NeuronData} c\^ot\'e Unity.

    \item \textbf{Stabilit\'e du mod\`ele en mode \'evaluation lors du
    backward pass} : l'utilisation de \texttt{BatchNorm} sur un batch de
    taille~1 peut poser probl\`eme.\\
    \textit{Solution} : conservation du mode \texttt{eval()} du r\'eseau,
    ce qui permet \`a \texttt{BatchNorm} d'utiliser ses statistiques
    courantes (\textit{running stats}) tout en laissant les gradients se
    propager normalement pour la mise \`a jour des poids.

    \item \textbf{Mise en place de l'environnement de r\'ealit\'e
    augment\'ee} : configuration de AR Foundation, gestion multi-plateforme
    (Android/iOS) et tests sur appareil physique, n\'ecessaires pour valider
    le placement et le suivi de la sc\`ene en conditions r\'eelles.
\end{itemize}


% =======================================================================
\chapter{Conclusion et perspectives}
\label{chap:conclusion}

\section{Conclusion}

Le projet \textbf{NeurAR} a permis de concevoir une cha\^ine compl\`ete
allant de l'entra\^inement d'un mod\`ele de r\'eseau de neurones avec
PyTorch, jusqu'\`a sa visualisation interactive en r\'ealit\'e augment\'ee
via Unity, en passant par un backend temps r\'eel bas\'e sur FastAPI et
WebSocket. Il r\'epond \`a une probl\'ematique p\'edagogique concr\`ete : la
difficult\'e de comprendre le fonctionnement interne des r\'eseaux de
neurones, en proposant une exp\'erience immersive, visuelle et interactive.

Les principaux r\'esultats obtenus sont :
\begin{itemize}
    \item un mod\`ele MLP fonctionnel, configurable et entra\^in\'e sur MNIST ;
    \item un backend robuste exposant l'\'etat du r\'eseau via REST et
    WebSocket, avec une animation couche par couche ;
    \item une application Unity AR capable d'afficher et d'animer le r\'eseau
    en 3D, avec des interactions utilisateur (rotation, zoom, s\'election).
\end{itemize}

\section{Perspectives d'am\'elioration}

Plusieurs pistes permettraient d'enrichir le projet :

\begin{itemize}
    \item \textbf{Support d'autres architectures} de r\'eseaux de neurones
    (r\'eseaux convolutifs CNN, r\'eseaux r\'ecurrents RNN) ;
    \item \textbf{Mode multi-utilisateurs} permettant \`a plusieurs personnes
    de visualiser et manipuler la m\^eme sc\`ene AR simultan\'ement ;
    \item \textbf{Entra\^inement interactif} : permettre de modifier les
    hyperparam\`etres (taux d'apprentissage, architecture) directement depuis
    l'application et observer l'impact en direct ;
    \item \textbf{Optimisation de la taille des donn\'ees transmises}, par
    exemple via la compression ou l'envoi de mises \`a jour diff\'erentielles
    (\textit{delta-update}) plut\^ot que l'\'etat complet \`a chaque \'etape ;
    \item \textbf{D\'eploiement sur casques AR/VR autonomes} (HoloLens,
    Meta Quest) pour une exp\'erience totalement immersive et mains libres.
\end{itemize}

\end{document}
